\chapter*{Resumen Ejecutivo}
\addcontentsline{toc}{chapter}{Resumen Ejecutivo}

El presente documento expone los resultados de una investigación forense digital exhaustiva realizada sobre un servidor departamental de producción (\texttt{victima1.foo.bar.mx}) perteneciente a una institución académica, tras el descubrimiento de actividades no autorizadas e implantación de software malicioso con privilegios administrativos.

La investigación abarcó dos fases metodológicas complementarias:
\begin{enumerate}
    \item \textbf{Análisis en vivo y monitoreo activo (\textit{live response}):} Ejecutado con el sistema en operación para contener riesgos inminentes, registrar procesos ocultos en memoria, interceptar comunicaciones de red y supervisar interactivamente las sesiones de cuentas sospechosas sin alertar prematuramente a los atacantes.
    \item \textbf{Análisis forense \textit{post-mortem}:} Conducción de exámenes minuciosos sobre imágenes binarias bit a bit de los discos duros en estaciones de trabajo de control y referencia (\textit{Dell Precision 350}). Esta fase comprendió el análisis de integridad de paquetes del sistema operativo (RPM), escaneo heurístico de malware, reconstrucción cronológica mediante líneas de tiempo de inodos (\textit{timeline analysis} con tiempos MAC), minería de bitácoras del sistema (\texttt{syslog}, \texttt{wtmp}, \texttt{xferlog}), recuperación de archivos eliminados (\textit{file carving}) y correlación multisistema con servidores de correo y autenticación web (\textit{OpenWebMail}).
\end{enumerate}

\begin{notaforense}[Hallazgo Principal de la Investigación]
La evidencia pericial recabada demostró concluyentemente la convergencia de \textbf{dos intrusiones independientes}:
\begin{itemize}
    \item Un \textbf{ataque externo oportunista}, originado a través de credenciales débiles en el servicio FTP, cuyo propósito principal consistió en utilizar la infraestructura del servidor como plataforma parásita para alojar servicios IRC, bots de relevo y depósitos de archivos.
    \item Un \textbf{segundo ataque de carácter interno}, iniciado de manera concomitante o inmediatamente posterior al primero, perpetrado por un actor con acceso a las instalaciones físicas del edificio y conocimiento de la topología local. Este segundo actor aprovechó el desconcierto causado por la primera intrusión para implantar herramientas de espionaje (\textit{sniffers}), sustituir binarios del sistema por \textit{rootkits} y capturar las nuevas contraseñas de administración de \texttt{root}.
\end{itemize}
\end{notaforense}

Las conclusiones del informe pericial permitieron deslindar responsabilidades, aislar los equipos vulnerados, erradicar la persistencia del malware en la red institucional e implementar políticas rigurosas de autenticación y segmentación de servicios.
